\immediate\write18{makeindex -s nomencl.ist -o Outputs/rapport.nls Outputs/rapport.nlo}

\documentclass[hidelinks,11pt]{article}

%Langue et encodage :
 \usepackage[T1]{fontenc}
 \usepackage[utf8]{inputenc}
 \usepackage{lmodern}
 \usepackage[french]{babel} \usepackage{lmodern}


%Mise en page
\usepackage[a4paper, margin=2.5cm]{geometry}	%Papier
\usepackage{arev}									%Police
\renewcommand{\baselinestretch}{1.15} 				%Interligne
\setlength\parindent{4pt}							%Alineas
\usepackage[skip=5pt,font=footnotesize]{caption}
\usepackage{float}
\usepackage{xcolor}									%Couleurs perso
\definecolor{grisPerso}{rgb}{0.9,0.9,0.9}
\usepackage{tabu}

%Contenu :
\usepackage{hyperref}								%Liens
\usepackage{amsmath}								%Math
\usepackage{listings}								%Code source
\lstset{
language=bash,
backgroundcolor=\color{grisPerso},
basicstyle=\footnotesize,
frame=single}
\usepackage[refpage]{nomencl}						%Nomenclature
\usepackage{xpatch}
\makenomenclature

%Dessins et schémas
\usepackage{graphicx}
\usepackage{epstopdf}
\epstopdfsetup{outdir=./}
\usepackage{subcaption}
\usepackage{tikz}
\usepackage{includes/tikz-uml}
\usetikzlibrary{calc,trees,positioning,arrows,chains,shapes.geometric,%
    decorations.pathreplacing,decorations.pathmorphing,shapes,%
    matrix,shapes.symbols,fit}


% Define block styles  
\tikzstyle{materia}=[draw, fill=blue!20, text width=6.0em, text centered,
  minimum height=1.5em,drop shadow]
\tikzstyle{module} = [materia, text width=8em, minimum width=10em,
  minimum height=3em, rounded corners, drop shadow]
\tikzstyle{linepart} = [draw, thick, color=black!50, -latex', dashed]
\tikzstyle{line} = [draw, thick, color=black!50]

\tikzstyle{myArrow} = [line, -latex']
\tikzstyle{ur}=[draw, text centered, minimum height=0.01em]

% Define distances for bordering
\newcommand{\blockdist}{1.3}
\newcommand{\edgedist}{1.5}

\newcommand{\module}[2]{node (p#1) [module]
  {#2}}
\include{includes/functions}
\begin{document}

%!TEX root = ../rapport.tex

\begin{titlepage}
\newgeometry{top=2cm,bottom=2cm}

\newcommand{\HRule}{\rule{\linewidth}{0.5mm}} % Defines a new command for the horizontal lines, change thickness here

\center % Center everything on the page
 
%----------------------------------------------------------------------------------------
%   HEADING SECTIONS
%----------------------------------------------------------------------------------------

\includegraphics[scale=0.6]{Graphics/LOGO-IPSA.pdf}\\[1.5cm]

\textsc{\Large Projet Electif}\\[0.5cm] % Major heading such as course name
\textsc{\large Projet technique}\\[0.5cm] % Minor heading such as course title

{\normalsize \today}\\[1.5cm]

%----------------------------------------------------------------------------------------
%   TITLE SECTION
%----------------------------------------------------------------------------------------

\HRule \\[0.4cm]
{ \Large \bfseries Evolution de véhicules autonomes dans un environnement urbain intelligent}\\[0.4cm] % Title of your document
\HRule \\[2.5cm]


%----------------------------------------------------------------------------------------
%   LOGO SECTION
%----------------------------------------------------------------------------------------

 \includegraphics[scale=0.4]{Graphics/illustration2.png}\\[2.5cm]

%----------------------------------------------------------------------------------------
%   AUTHOR SECTION
%----------------------------------------------------------------------------------------
{\normalsize Auteurs:}\\
\small
\textsc{Biton} Guillaume (guillaume.biton@ipsa.fr)\\
\textsc{Guichard} Marc-Antoine \small(marc-antoine.guichard@ipsa.fr)\\
\textsc{Lhermite} Camille \small(camille.lhermite@ipsa.fr)\\
\textsc{Monnot} Maxime \small(maxime.monnot@ipsa.fr)\\[1cm]

 
%----------------------------------------------------------------------------------------

\vfill % Fill the rest of the page with whitespace

\restoregeometry

\end{titlepage}

\setcounter{secnumdepth}{5}
\setcounter{tocdepth}{5}

\newpage
\tableofcontents
\newpage
\pagenumbering{arabic}

\section*{Remerciements}

	Nous remercions Mesdames Assia Belbachir et Sorore Benabid, ensignants chercheurs au sein du laboratoire Mécatronique, Signal et Systèmes de l'IPSA pour nous avoir soumis cette idée, pour nous avoir accordé leur confiance quant à sa réalisation potentielle et pour leur soutien dans le cadre de ce projet.

	\vspace{80pt}

\section*{Notes préliminaires}
	L'ensemble des illustrations présentes dans ce dossier, sauf indication contraire dans leurs légendes, sont des réalisations personnelles et ne sont donc soumises à aucun droit particulier.\\

	Le travail de recherche ici présenté est sous la responsabilité des auteurs cités en première de couverture et sous la propriété de l'IPSA. Merci de nous contacter au moyen des adresses fournies pour nous signaler tout problème, ou pour toute demande liée à l'utilisation du contenu.\\

	L'ensemble des abréviations utilisées dans ces pages sont explicitées en page \pageref{nomenclature} (nomenclature).\\

	Les numéros entre crochets pouvant apparaitre au sein de paragraphes sont des références bibliographiques et vous renvoient à la page \pageref{bibliographie} (bibliographie).\\

	Ce dossier comprend de nombreuses références dynamiques, et pour un confort de lecture optimal, nous vous conseillons d'en lire la version numérique afin de pouvoir bénéficier des liens cliquables.\\

	Exception faite des modèles mécaniques en CAO 3D effectués sous Catia\textsuperscript{\textregistered}, l'ensemble des réalisations ici présentées utilisent des outils Open-Source : le présent dossier est rédigé en \href{https://www.latex-project.org/}{\LaTeX}, les illustrations ont été réalisées en dessin-vectoriel sous \href{https://inkscape.org/fr/}{Inkscape} ou en matriciel sous \href{https://www.gimp.org/}{GIMP} et les schémas et circuits électroniques sous \href{http://kicad-pcb.org/}{KiCad}.

\newpage

\section{Introduction}

	La voiture autonome, loin du concept de science-fiction qu'elle pouvait représenter il y a quelques années est en train de devenir une réalité.\\

Si cette ambition put être à une certaine époque motivée par le simple attrait de la prouesse technique, nous percevons aujourd'hui tous les bénéfices que l'on pourrait en tirer.
En effet, les avancées scientifiques et techniques nous permettent désormais de prétendre à concevoir une voiture qui soit non seulement autonome, mais surtout intelligente. Il est aujourd'hui tout à fait réaliste de penser que dans les quelques années à venir les voitures sauront adopter un comportement bien plus intelligent que celui de leurs conducteurs actuels, et ce au profit de la sécurité, de l' efficience énergétique mais également de l'encombrement des axes routiers.\\

A terme, nous pouvons facilement imaginer que les différents véhicules auront la possibilité de communiquer entre eux afin de prévenir les véhicules environnants de leurs intentions, mais cela ne les dispensera pas de devoir être capables d'évaluer leur environnement afin d'y détecter les éléments "indépendants" (piétons, obstacles...). \\

Comme toute révolution technologique, la voiture intelligente devra faire face au caractère progressif de son adoption: toutes les voitures sur les routes ne deviendront pas autonomes du jour au lendemain. Ces véhicules devront donc également être capables d'évoluer au milieu d'une circulation telle que nous la connaissons, où chaque acteur adopte un comportement presque parfaitement aléatoire, et ne signale pas toujours ces intentions.\\

Si les voitures deviennent intelligentes, c'est également le cas, depuis quelques années déjà, de la signalisation. Ainsi, de plus en plus de feux adaptent leur comportement au trafic, ce qui s'inscrit également dans une démarche d'optimisation de la circulation. Il est très réaliste d'espérer que cette technologie déjà existante se fera omniprésente dans les années qui viennent, d'où notre volonté d'inclure cet élément d’environnement à notre projet.\\

Le but de ce projet est donc d'étudier notre capacité à faire cohabiter intelligence artificielle et environnement "réel" et indépendant avec des moyens techniques et financiers extrêmement restreints, mais également et surtout de fournir une plateforme d'expérimentation aux étudiants de l'IPSA (et d'ailleurs ?) afin de faciliter l'apprentissage des sciences de la mécatronique et de l’intelligence artificielle tout en ancrant ce savoir dans une réalité pratique et stimulante.

\newpage
\section{Définition du besoin}

	Une définition pertinente du besoin est une condition absolument nécessaire pour la bonne réalisation de tout projet. Nous nous emploierons donc à définir aussi précisément et pertinemment que possible le besoin motivant ce projet, et à régulièrement revenir sur ce dernier afin de prendre en compte d'éventuelles évolutions et à prévenir toute dérive du projet.

\subsection{Besoin}
Le besoin à l'origine de ce projet fut exprimé par mesdames A. Belbachir et S. Benabid en tant qu'enseignants chercheurs au laboratoire Mécatronique, Signal et Systèmes à l'IPSA. Leur souhait était de bénéficier d'une plateforme articulée autour de robots autonomes évoluant dans une simulation d’environnement urbain, capables de détecter des feux de signalisation "intelligents" et d'adapter leur comportement en conséquence. Cette plateforme pourrait servir de support de TP\nomenclature{TP}{Travaux Pratiques} dans l'ensemble des matières enseignées par le département, mais également servir de "vitrine" voir même de plateforme de recherche.\\

Les principales fonctionnalités évoquées étant :
\begin{itemize}
	\item \textbf{La reconnaissance d'image} pour la détection des feux de signalisation.
	\item \textbf{La présence de feux bicolores commandés par FPGA\nomenclature{FPGA}{PGA (Field-Programmable Gate Array, ou Réseau de Portes Programmables : un type très répandu de circuit logiques programmables.}.} Ces feux seraient "intelligents" dans le sens où il s'adapteraient à la circulation. La présence de capteurs de présence est donc induite.
\end{itemize}

\subsection{Précision du besoin}

Ayant suivi nombre de matières enseignées par le département et participé à de nombreuses séances de travaux pratiques, nous bénéficions d'une image relativement claire des implications de ce projet ainsi que des contraintes relevant pour la plupart du simple bon sens.
Dans un soucis de clarté et d'application de "bonnes pratiques" nous prîmes cependant soin de valider l'ensemble de ces éléments au cours de réunions avec nos "clients" du laboratoire.\\

Ainsi, nous ajoutâmes les points suivants à la liste des exigences :
\begin{itemize}
	\item "Côté circuit" :
	\begin{itemize}
		\item \textbf{Le circuit devra bénéficier d'un encombrement raisonnable}.
		\item \textbf{Les feux de circulation devront pouvoir être commandés via une carte FPGA}.
	\end{itemize}
	\item "Côté robots" :
	\begin{itemize}
		\item \textbf{Les robots devront pouvoir être utilisés en classe sans que les préoccupations matérielles ne soient accaparantes}.
		\item \textbf{Les robots devront pouvoir être programmés en utilisant les langages enseignés à l'IPSA} à savoir C et C++ et ses variantes (Arduino...), Python, voir même Matlab...
		\item \textbf{Les robots devront bénéficier de possibilités d'application flexibles :} les enseignements étant ciblés, il est important que les utilisateurs des robots puissent se concentrer sur un aspect de leur utilisation sans avoir à se soucier des autres. De même, les robots devront embarquer suffisamment de capteurs ou tout du moins de capacité d'extensions pour que cela ne soit pas un facteur limitant lors de l'élaboration de sujets de TP.
		\item \textbf{Les robots devront ne pas pouvoir représenter un danger pour ses utilisateurs}.
	\end{itemize}
\end{itemize}
\newpage
\section{Étude du besoin}

	\subsection{Analyse du besoin}

Le besoin exprimé précédemment comporte un certain nombre d'implications logiques.
Leur prise en compte relève déjà du domaine de la solution, mais reste d'ordre suffisamment général pour être évoquée ici. Surtout, ces points seront à prendre en compte lors de l'élaboration d'une solution "détaillée".\\

Reprenons les points évoqués ci-avant:

\renewcommand{\labelitemi}{\textbullet}
\renewcommand{\labelitemii}{$\Rightarrow$}
\renewcommand{\labelitemiii}{-}
\renewcommand{\labelitemiv}{\textbullet}

	\subsubsection{Pour le circuit}
		\begin{itemize}
			\item \textbf{Le circuit devra bénéficier d'un encombrement raisonnable}.
			\begin{itemize}
				\item Les dimensions du circuit devront être raisonnables.
				\item Nous songerons également à une structure démontable (qui devra rester simple à assembler).
			\end{itemize}
			\item \textbf{Les feux de circulation devront pouvoir être commandés via une carte FPGA}.
			\begin{itemize}
				\item Une carte adaptant les tensions et courants et disposant de connecteurs adaptés aux caractéristiques des cartes possédées par l'IPSA devra être proposée.
				\item La carte FPGA aura pour rôle de synchroniser les feux.
				\item La carte FPGA pourra tirer partie de capteurs de présence par adapter le cycle des feux.
			\end{itemize}
		\end{itemize}
	\subsubsection{Pour le robot}
		\begin{itemize}
			\item \textbf{Les robots devront pouvoir être utilisés en classe sans que les préoccupations matérielles ne soient accaparantes.}
			\begin{itemize}
				\item Cela implique plusieur points :
				\begin{itemize}
					\item Bénéficier d'une autonomie suffisante pour que l'ensemble des groupes puissent effectuer leurs manipulations au cours d'une séance entière sans que cela ne nécessite de charger/remplacer la batterie. D'expérience, nous estimons l'autonomie minimum à une heure pour répondre à ce critère.
					\item Présenter des connecteurs et des interfaces permettant une interaction "technique" avec le robot ne nécessitant pas de matériel complexe. Ceci passe par la possibilité de changer ou recharger la batterie sans outils particuliers, de pouvoir programmer le robot via USB ou réseau plutôt que via un port parallèle par exemple, ou encore proposer au minimum un bouton accessible permettant le lancement d'un séquence de code prédéfinie.
					\item Etre d'une conception suffisamment robuste pour pouvoir être manipulé quotidiennement sans en "souffrir" et suffisamment simple pour pouvoir être aisément entretenu.
				\end{itemize} 
			\end{itemize}
			\item \textbf{Les robots devront pouvoir être programmés en utilisant les langages enseignés à l'IPSA.}
			\begin{itemize}
				\item Cette problématique sera étudiée plus en détail en \ref{Solution-controleur} (page \pageref{Solution-controleur})
			\end{itemize}
			\item \textbf{Les robots devront bénéficier de possibilités d'application flexibles.}
			\begin{itemize}
				\item Ceci passe par une \textbf{modularité totale} des différents éléments qui composent le robot, aussi bien sur le plan matériel que logiciel. Cette modularité devra rester une préoccupation de premier plan tout au long du développement de ce projet.
			\end{itemize}
			\item \textbf{Les robots devront ne pas représenter de danger pour ses utilisateurs}.
			\begin{itemize}
				\item La structure du robot ne devra pas présenter d'éléments tranchants.
				\item La partie électrique et électronique devra employer des puissances raisonnables, et être protégée contre les forts courants.
			\end{itemize}
		\end{itemize}


\subsection{Diagrammes des cas d'utilisation}

Toujours dans un soucis de clarté et de bonne pratiques, nous avons ébauché les diagrammes des cas d'utilisation de haut niveau suivants :

	\subsubsection{Pour le circuit}

		\begin{tikzpicture} 
			\setcounter{tikzumlUseCaseNum}{0}
			\begin{umlsystem}[x=10, fill=blue!10]{Circuit} 
				\umlusecase[fill=white]{Utilisation expérimentale} 
				\umlusecase[y=-2,fill=white]{Utilisation "démonstration"} 
				\umlusecase[y=-4,fill=white]{Mettre en place le circuit}
				\umlusecase[y=-5.5,fill=white]{Ranger le circuit}  
			\end{umlsystem} 
			 
			\umlactor[y=0.5]{Développeur}  
			\umlactor[y=-2.5]{Utilisateur} 
			   
			\umlinherit{Développeur}{Utilisateur} 
			\umlassoc{Développeur}{usecase-1} 
			\umlassoc{Utilisateur}{usecase-2} 
			\umlassoc{Utilisateur}{usecase-3}  
			\umlassoc{Utilisateur}{usecase-4}  
		\end{tikzpicture}

		\textbf{Utilisation expérimentale}
		\begin{itemize}
			\item \textbf{Précondition :} Le circuit est en place. L'utilisateur dispose d'une carte FPGA programmée et alimentée.
			\item \textbf{Déclencheur :} Un utilisateur souhaite tester son programme VHDL
			\item \textbf{Scenario :}
			\begin{enumerate}
				\item L'utilisateur connecte sa carte FPGA à la carte du circuit en veillant à respecter les branchements.
				\item L'utilisateur effectue les tests et observations qu'il désire.
				\item L'utilisateur déconnecte sa carte FPGA.
			\end{enumerate}
		\end{itemize}
		\newpage

		\textbf{Utilisation "démonstration"}
		\begin{itemize}
			\item \textbf{Précondition :} On dispose d'une carte programmée pour la demonstration.
			\item \textbf{Déclencheur :} Un utilisateur souhaite effectuer une démonstration.
			\item \textbf{Scenario :}
			\begin{enumerate}
				\item L'utilisateur connecte la carte FPGA de démonstration à la carte du circuit en veillant à respecter les branchements.
				\item L'utilisateur alimente la carte FPGA de démonstration.
				\item L'utilisateur laisser le système évoluer en autonomie.
				\item L'utilisateur coupe l'alimentation de la carte FPGA de démonstration.
				\item L'utilisateur déconnecte laa carte FPGA de démonstration.
			\end{enumerate}
		\end{itemize}

		\textbf{Mettre en place le circuit}
		\begin{itemize}
			\item \textbf{Précondition :} On dispose es éléments du circuit et de la place nécéssaire à son installation
			\item \textbf{Déclencheur :} Un souhaite utiliser le circuit qui n'est pas déjà en place
			\item \textbf{Scenario :}
			\begin{enumerate}
				\item L'utilisateur s'assure que les deux parties sont en état d'être assemblées.
				\item L'utilisateur isole les deux faisceaux de câbles.
				\item L'utilisateur assemble les deux parties du circuit.
				\item L'utilisateur verrouille les deux parties du circuit.
				\item L'utilisateur installe les feux à leurs emplacements et les connecte à leurs faisceaux de câbles.
				\item L'utilisateur installe les différents "caches".
				\item L'utilisateur branche les deux faisceaux de câbles à leurs connecteurs sur la carte du circuit.
			\end{enumerate}
		\end{itemize}

		\textbf{Ranger le circuit}
		\begin{itemize}
			\item \textbf{Précondition :} Le circuit est installé et n'est pas en cours d'utilisation.
			\item \textbf{Déclencheur :} Un ne se sert pas du circuit et on souhaite libérer la place.
			\item \textbf{Scenario :}
			\begin{enumerate}
				\item L'utilisateur déconnecte les deux faisceaux de câbles de la carte du circuit.
				\item L'utilisateur retire les différents "caches".
				\item L'utilisateur déconnecte les feux de leurs faisceaux de câbles et les retire de leurs emplacements.
				\item L'utilisateur déverrouille les deux parties du circuit.
				\item L'utilisateur s'assure que les deux parties sont en état d'être manipulées et déplacées.
				\item L'utilisateur range les deux parties du circuit.
			\end{enumerate}
		\end{itemize}

	\subsubsection{Pour le robot}

		\begin{tikzpicture} 
			\begin{umlsystem}[x=10, fill=blue!10]{Robot} 
			\umlusecase[fill=white]{Charger un programme} 
			\umlusecase[y=-2,fill=white]{Evoluer en autonomie} 
			\umlusecase[y=-4,fill=white]{Changer la batterie} 
			\end{umlsystem} 
			 
			\umlactor{Développeur}  
			\umlactor[y=-3]{Utilisateur} 
			   
			\umlinherit{Développeur}{Utilisateur} 
			\umlassoc{Développeur}{usecase-1} 
			\umlassoc{Utilisateur}{usecase-2} 
			\umlassoc{Utilisateur}{usecase-3}	 
		\end{tikzpicture} 

		\textbf{Reprogrammation :}
		\begin{itemize}
			\item \textbf{Précondition :} Nous disposons d'un programme fonctionnel.
			\item \textbf{Déclencheur :} Un développeur souhaite implémenter un programme.
			\item \textbf{Scenario :}
			\begin{enumerate}
				\item On connecte le robot à une source de tension.
				\item On établit une liaison entre le robot et un ordinateur.
				\item Le développeur lance le téléversement du programme.
				\item Le robot acquitte.
				\item On ferme la connexion.
			\end{enumerate}
		\end{itemize}

		\textbf{Evolution en autonomie:}
		\begin{itemize}
			\item \textbf{Précondition :} Le robot a été programmé et dispose d'une batterie chargée.
			\item \textbf{Déclencheur :} Un membre du laboratoire souhaite observer le comportement du robot.
			\item \textbf{Scenario :}
			\begin{enumerate}
				\item On met le robot sous tension.
				\item On place le robot sur une ligne blanche du circuit.
				\item On appuie sur le bouton de démarrage de séquence
				\item Le robot déclenche les programmes en mémoire.
			\end{enumerate}
		\end{itemize}

		\textbf{Changement de batterie:}
		\begin{itemize}
			\item \textbf{Déclencheur :} On souhaite changer la batterie du robot
			\item \textbf{Scenario :}
			\begin{enumerate}
				\item On met le robot hors tension.
				\item On accède à la batterie en place le cas échéant, et on la retire.
				\item On met en place la nouvelle batterie
				\item On remet en place les éléments éventuellement retirés pour accéder à la batterie.
			\end{enumerate}
		\end{itemize}


		Le mode d'utilisation primaire est bien évidemment celui de l'évolution en autonomie. Le robot étant destiné à servir de plateforme de recherche, et devant donc être entièrement reprogrammable, il est délicat de décrire ce mode d'utilisation qui dépendra intégralement du programme chargé par l'utilisateur.\\

		Nous nous appliquerons cependant à décrire le mode d'utilisation correspondant à l'application la plus basique du robot (et de fait au programme que nous livrerons avec ce dernier).
		

\newpage
\section{Élaboration d'une solution}

	\input{chapitres/3_elaborationSolution}

\newpage
\section{Intégration}

	\input{chapitres/4_integration}

\newpage
\section{Estimation du coût de mise en place et faisabilité}

	\input{chapitres/5_coutFaisabilite}

\section{Suivi de projet}

	Un suivi de projet rigoureux devra être mis en place.\\

Nous recommandons l'instauration d'un compte-rendu hebdomadaire (qui pourra être fait par mail), ainsi que d'une réunion d'avancement mensuelle.\\

Il serait également judicieux de se contraindre à rédiger quotidiennement un journal de travail.\\

Surtout, un tableau partagé permettra de suivre l'avancement "en temps réel" et, si besoin, de réajuster les prévisions et l'allocation des ressources.\\

Nous avons produit le "code de bonne pratiques" suivant afin d'avoir une référence à laquelle se rapporter au cours du projet:\\
\begin{adjustwidth}{1cm}{0.5cm}
	\fbox{\parbox{\linewidth-2\fboxrule-2\fboxsep}{
	Ce guide de bonnes pratiques est à consulter régulièrement et respecter tout au long du projet.\\
	Ce guide a vocation à améliorer la productivité et non l'inverse : les protocoles "lourds" et présentant peu d'intérêt sont à proscrire.\\

	\begin{enumerate}
		\item \textbf{COMMUNICATION :}
		\begin{itemize}
			\item \textbf{EMAILS :}\\
			Pour informations qui doivent rester, et/ou échange avec différents acteurs\\
			\textbf{Objet :} [PFE VA] - [CATEGORIE] Objet	\textit{(ex: [PFE VA] - [FINANCIER] Dépassement de budget)}\\
			\textbf{Contenu :} Concis, en "points par points"\\
			Utiliser les champs "Destinataire, "Copie Carbone" et "Copie Carbone Invisible" à bon escient (Destinataire pour personne directement visée, CC pour acteurs à tenir informés)\\
			Essayer de répondre en point par point sur la structure du mail initial\\

			\item \textbf{SLACK :}\\
			Pour communication "temps réel", conversation. Pas de décisions importantes.\\
			Diviser en canaux par thématique\\
		\end{itemize}

		\item \textbf{STOCKAGE :}\\
		Respecter l'architecture en place\\
		Noms de fichiers doivent être clairs, explicites et concis\\
		Sauvegardes régulières sur support externe\\
	\end{enumerate}}}
\end{adjustwidth}
\begin{adjustwidth}{1cm}{0.5cm}
	\fbox{\parbox{\linewidth-2\fboxrule-2\fboxsep}{
	\begin{enumerate}
	\setcounter{enumi}{2}
		\item \textbf{GIT :}\\
		Toujours commencer session de travail par un pull\\
		Un commit à chaque "bloc de modification" : ajout d'une fonction, correction d'un bug... Pas de commit sur une modif non testée\\
		Création d'une branche pour chaque modification "expérimentale". Fusion de la branche si concluante, après concertation.\\
		Exclure des commits tous les fichiers binaires (sauf illustrations quand utilisées par le code) et tous les fichiers générés (fichiers de compilation...) : git ne doit concerner que les fichiers textes (sauf fichiers nécessaires à la bonne compilation/utilisation du code). => utiliser le .gitignore\\
		=> un pull doit toujours être suivi d'une recompilation : ne jamais  utiliser un fichier déjà compilé (qui ne doivent de toute façon pas être commit).

		
		\item \textbf{CODE :}
		Respect de la convention décrite ici : \url{http://all4dev.libre-entreprise.org/index.php/Conventions_de_syntaxe_en_python}
		Respect de la convention PEP8
		Ajout systématique d'un docstring avec description de la fonction (ou classe), des paramètres et du retour.

		\item \textbf{SUIVI :}
		Utilisation de ProjectLibre pour ajuster le planning et l'avancement en temps réel
		Un point d'avancement par email 1 fois par semaine avec accomplissements, problèmes rencontrés, besoins éventuels et modifications notables du projet.

		\item \textbf{DOCUMENTATION :}
		\begin{itemize}
			\item \textbf{EXISTANTE :}\\
			Prendre soin de conserver toute la documentation dans le répertoire prévu à cet effet. Veiller à suivre une nomenclature et une organisation simple et explicite.\\
			Compléter la bibliographie au fur et à mesure via le formulaire dédié.\\
			\item \textbf{PRODUITE :}\\
			Au fur et à mesure du développement, utiliser les commentaires et les docstrings pour documenter les programmes.\\
			Enregistrer dans de simples fichiers textes, ou communiquer par email s'ils doivent être partagés les points à prendre en compte \textit{(ex: Attention à toujours couper le contact avant de charger un programme !)}\\
			La construction de la documentation livrable ne débutera qu'en fin de projet, quand nous bénéficierons d'une vue d'ensemble.\\
		\end{itemize}
	\end{enumerate}}}
\end{adjustwidth}
\newpage
\section{Conclusion et aperçu des évolutions possibles}
	
	\input{chapitres/7_conclusion}

\setcounter{secnumdepth}{0}
\renewcommand{\thesubsection}{\Alph{subsection}}

\newpage
\sectionWoNum{Annexes}
	\subsection{Carte Circuit}\label{schemasCarteMere}
	\subsubsection{Schéma électronique}

		\illu{ANNEXES/SCHEMA_CARTE_CIRCUIT.pdf}{Schéma du circuit "circuit"}{0.8}

		\newpage
	\subsubsection{Gerbers}
		\illu{ANNEXES/CARTE_CIRCUIT.pdf}{Gerbers de la carte "circuit"}{1}
	\subsection{Carte Réflecteurs Optiques}\label{schemasCarteReflecteurs}
	\subsubsection{Schéma électronique}

		\illu{ANNEXES/CARTE-REFLECTEURS_circuit.pdf}{Schéma électronique de la carte Réflecteurs Optiques}{1}

		\newpage
	\subsubsection{Gerbers}
		\illu{ANNEXES/CARTE-REFLECTEURS.pdf}{Gerbers de la carte Réflecteurs Optiques}{1}
	\subsection{Carte Mère}\label{schemasCarteMere}
	\subsubsection{Schéma électronique}

		\illu{ANNEXES/SCHEMA_CARTE_MERE_1.pdf}{Schéma du circuit de la carte mère}{1}
		\illu{ANNEXES/SCHEMA_CARTE_MERE_2.pdf}{Schéma du circuit de la carte mère}{1}


		\newpage
	\subsubsection{Gerbers}
		\illu{ANNEXES/CARTE_CARTE_MERE.pdf}{Gerbers de la carte mère}{1}
	\newpage
	\subsection{Dessins Techniques structure}\label{devis}
	\subsubsection{Chassis arrière}

		\illu{ANNEXES/STR_chassis_ar_01.pdf}{Structure : chassis arrière}{0.52}

		\newpage

	\subsubsection{Chassis avant}

		\illu{ANNEXES/STR_chassis_av_01.pdf}{Structure : chassis avant}{0.52}

	\subsubsection{Chassis droit}

		\illu{ANNEXES/STR_chassis_dr_01.pdf}{Structure : chassis droit}{0.52}

		\newpage

	\subsubsection{Chassis gauche}

		\illu{ANNEXES/STR_chassis_ga_01.pdf}{Structure : chassis gauche}{0.52}

		\newpage

	\subsubsection{Fond}

		\illu{ANNEXES/STR_fond_01.pdf}{Structure : Fond}{0.52}

		\newpage
	\subsubsection{Montant arrière-droit}

		\illu{ANNEXES/STR_mont_ar-dr_01.pdf}{Structure : Montant arrière-droit}{0.52}

		\newpage
	\subsubsection{Montant arrière-gauche}

		\illu{ANNEXES/STR_mont_ar-ga_01.pdf}{Structure : Montant arrière-gauche}{0.52}

		\newpage
	\subsubsection{Montant avant-droit}

		\illu{ANNEXES/STR_mont_av-dr_01.pdf}{Structure : Montant avant-droit}{0.52}

		\newpage
	\subsubsection{Montant avant-gauche}

		\illu{ANNEXES/STR_mont_av-ga_01.pdf}{Structure : Montant avant-gauche}{0.52}

		\newpage
	\subsubsection{Support Webcam}

		\illu{ANNEXES/STR_support-cam_01.pdf}{Structure : Support Webcam}{0.52}

		\newpage

	\subsubsection{Support Réflecteurs}

		\illu{ANNEXES/STR_support-reflecteurs_01.pdf}{Structure : Support Reflecteurs}{0.52}

		\newpage
	\subsection{Devis matériel}\label{devis}
	\subsubsection{Devis global}

		\illu{ANNEXES/DEVIS_global.pdf}{Devis global}{1}

		\newpage

	\subsubsection{Devis pour un robot}
		\illu{ANNEXES/DEVIS_robot.pdf}{Devis pour un robot}{1}

		\newpage

		\paragraph{Devis pour la carte-mère}\mbox{} \\

			\illu{ANNEXES/DEVIS_CARTE_mere.pdf}{Devis pour une carte-mère}{1}

		\paragraph{Devis pour la care réflecteurs}\mbox{} \\
			\illu{ANNEXES/DEVIS_CARTE_reflecteurs.pdf}{Devis pour la carte réflecteurs}{1}

	\subsubsection{Devis pour le circuit}
		\illu{ANNEXES/DEVIS_circuit.pdf}{}{1}

		\paragraph{Devis pour la carte du circuit}\mbox{} \\

			\illu{ANNEXES/DEVIS_CARTE_circuit.pdf}{Devis pour la carte du circuit}{1}
	\newpage
\subsection{Diagramme de Gantt}\label{gantt}
	
	\includegraphics[angle=90,origin=c,scale=0.45]{Graphics/Gantt.png}
	\newpage
\subsection{Diagramme de Gantt}\label{PERT}
	
	\includegraphics[angle=90,origin=c,scale=0.9]{Graphics/PERT1.png}
	\newpage
	\includegraphics[angle=90,origin=c,scale=0.6]{Graphics/PERT2.png}
\setcounter{secnumdepth}{0}
\newpage
\section{Références}

	\subsection{Nomenclature} \label{nomenclature}
		\printnomenclature

	\newpage
	\nocite{bib11,bib1}
	\subsection{Bibliographie} \label{bibliographie}
		\begingroup
		\renewcommand{\section}[2]{}%
		\bibliography{includes/bibliography}{}
		\bibliographystyle{plain}
		\endgroup

	\newpage
	\subsection{Table des illustrations}
		\begingroup
		\renewcommand{\section}[2]{}%
		\listoffigures
		\endgroup
		
\newpage
\section{Abstract}
	Ce dossier constitue l'\textbf{étude préliminaire} d'un projet de réalisation de \textbf{plateforme d'enseignement et de recherche} pour le pôle Mécatronique, Signal et Systèmes de la Direction de la Recherche et de l’Innovation de l’IPSA (DRII) de l'IPSA, école d'ingénieur de l'air et de l'espace.\\

Cette plateforme, composée d'un \textbf{"circuit" interactif} modelisant un réseau routier urbain et de \textbf{"robots" autonomes} capables de se déplacer en autonomie et en respect des règles de la route sur ce circuit, pourra servir de \textbf{support de travaux pratiques} pour l'ensemble des enseignements du pôle, voir même de \textbf{support de recherche}.\\

Le présent dossier propose une \textbf{analyse détaillée} de ce \textbf{besoin} et de ses implications, ainsi que la description d'une \textbf{solution technique} complète et argumentée. Enfin, une estimation des \textbf{coûts} et des \textbf{plannings} est offerte.\\

Ce dossier est agrémenté de diverses illustrations, schémas, plans et vues 3D, formant ainsi un guide presque suffisant pour la réalisation du projet.
\end{document}